\section{The survivor and its reversal}

Let $\Sigma_A$ be the edge shift on the alphabet $\{0,1,2,3\}$ with
\begin{equation}\label{eq:adjacency}
A=
\begin{pmatrix}
1&0&1&0\\
1&0&0&0\\
0&1&0&1\\
0&1&0&0
\end{pmatrix}.
\end{equation}
Write $\rho=(0)(1\ 2)(3)$ and let $P$ be its permutation matrix.

\begin{proposition}[Reversal identity]\label{prop:reversal}
The involution $\rho$ is the unique symbol permutation satisfying
\begin{equation}\label{eq:reversal}
A=P A^{\mathsf T}P.
\end{equation}
Consequently, for every cyclic admissible word $w=(w_j)_{j\in\mathbb Z/n}$,
\begin{equation}\label{eq:Rk}
(R_k w)_j=\rho(w_{-j-k})
\end{equation}
is again admissible, and $R_k^2=1$.
\end{proposition}

\begin{proof}
Direct multiplication proves \eqref{eq:reversal}.  Exhaustion of the
$4!$ symbol permutations proves uniqueness.  In entry form the identity is
$A_{ij}=A_{\rho(j),\rho(i)}$, which reverses every admissible edge; applying
$\rho$ twice proves the involution statement.
\end{proof}

The characteristic polynomial factors as
\begin{equation}\label{eq:charpoly}
\det(zI-A)=(z^2+1)(z^2-z-1).
\end{equation}
Let $F_m$ and $L_m$ denote the Fibonacci and Lucas sequences with
$F_0=0,F_1=1$ and $L_0=2,L_1=1$.  Taking powers in
\eqref{eq:charpoly} gives the closed-word identity
\begin{equation}\label{eq:closed}
N_n:=\Tr(A^n)=L_n+2\cos(\pi n/2).
\end{equation}

For even periods we name the two classes \emph{edge--edge} ($\EE$) and
\emph{vertex--vertex} ($\VV$) according to the physical H\'enon closure used
in the predecessor project.  These names should not be inferred merely from
which vertices of an abstract $n$-gon are fixed by \eqref{eq:Rk}.  The exact
word locks in Section~6 remove that possible convention ambiguity.
